\documentclass{exam}

\usepackage[margin=1in]{geometry}

\usepackage{comment}
\usepackage{graphicx}
\usepackage{listings}

\pagestyle{headandfoot}
\firstpageheader{CSCI 221}{}{Floating Point Assignment}
\runningheader{CSCI 221}{}{Floating Point Assignment}


\begin{document}
\title{Floating Point Assignment (Evaluation)}
\author{CSCI 221: Computer Science Fundamentals 2}
\date{Spring 2026}
\maketitle


This assignment is designed to evaluate your understanding of floating point numbers. 
Feel free to collaborate with others and use resources, but all code and writeups must be your own. 
To submit your assignment, please submit your code, makefile, and writeup. 
For full credit, your code should check for invalid input. 

\textbf{Due Date:}
Wednesday, March 18th at 1:00 pm. 

For this assignment you will be dealing with Tiny Floating Point numbers. 
These numbers use the same IEEE standard encoding schemes that we talked about in class. 
However, they contain only 8 bits. 
The MSB is the sign bit. 
The significand representation is the 3 least significant bits. 
The remaining 4 bits are the exponent representation. 
You will be representing Tiny Floating Point numbers as a sequence of bits using an appropriate type. 
For full credit, do not use an array, a struct, or any other composite data type in your representation. 

\begin{questions}
\question[10]
\textbf{Floating Point Increments.}
\begin{parts}
\part[5]
What is the difference between the largest non-infinite number representable with tiny floating point and the second largest such number?
Show your work. 

\textbf{Answer:}
Since there are 4 bits for encoding the exponent and 3 bits for the significand. The largest exponent encoding that is not all ones is 
$$1110_2=2^4-2=14.$$
The bias is $2^{4-1}-1=8-1=7$.
Thus, the largest exponent is
$$14-7=7.$$
The largest significand encoding is all ones, which means the largest significand has $3$ ones after the decimal point:
$$1.111=2-2^{-3}$$
which  This is
Thus, the largest number is
$$(2-2^{-3})\cdot2^{7}=2^{8}-2^{4}.$$
The second largest number have a $0$ instead of a $1$ in the least significant digit of the significand. The significand is
$$1.110=2-2^{-2}$$
Thus, it is
$$(2-2^{-2})\cdot2^{7}=2^{8}-2^{5}.$$
Finally, the difference is
$$2^{8}-2^{4}-(2^{8}-2^{5})=2^{5}-2^{4}=32-16=16.$$

\part[5]
What is the difference between the second smallest positive number representable with tiny floating point and the smallest such number?
Show your work. 
\end{parts}

\textbf{Answer:}
The smallest exponent encoding is $0$. Then its in the denormalized encoding. In this case, the smallest non-zero significand is all $0$s and a $1$ in the least significant digit. So the significand is $0.001=2^{-3}$ while the exponent value is $0-7+1=-6$. Thus, the value is
$$2^{-3}\cdot2^{-6}=2^{-9}.$$
And the next smallest is
$$2^{-2}\cdot2^{-6}=2^{-8}$$
The difference is
$$2^{-8}-2^{-9}=2^{-9.}$$

\question[20]
\textbf{Adding Tiny Floating Point.}
Write a function that takes two tiny floating point numbers as input and returns their sum as a tiny floating point number. 
Hint: You should be able to combine a significant amount of functionality between adding and multiplying. 

\question[20]
\textbf{Multiplying Tiny Floating Point.}
Write a function that takes two tiny floating point numbers as input and returns their product as a tiny floating point number. 
Hint: You should be able to combine a significant amount of functionality between adding and multiplying. 

\question[10]
\textbf{Writeup.}
In addition to your code, you should submit a brief (~1-2 pages) writeup that describes the state of your code. 
You should discuss what is working, what is not, and any known errors. 
In this discussion, you should reference the tests you did and how they support your claims. 
You will need to provide a brief description of your tests for you to reference. 
Your writeup should be in a separate .txt or .pdf file. 
\end{questions}

\end{document}